\documentclass[10pt]{article}
\usepackage[T1]{fontenc}
\usepackage[utf8]{inputenc}
\usepackage[french]{babel}
\usepackage{amsmath,amssymb,mathrsfs,mathpazo}
\usepackage{geometry}
\geometry{a4paper,margin=2cm}
\usepackage{graphicx,array,tabularx,tabularray,colortbl,multirow,multicol,booktabs,enumitem,pifont,ulem,fancybox,diagbox,cancel}
\usepackage{pstricks,pstricks-add,pst-plot,pst-tree,pst-node,pst-text,pst-func,pst-math,pst-bspline,pst-3dplot,pst-eucl,pst-solides3d}
\let\pscalculate\undefined
\usepackage{pst-bezier}
\usepackage{tikz,pgfplots,tkz-tab}
\usetikzlibrary{arrows,arrows.meta,patterns,calc,positioning,trees,decorations.markings,shapes.geometric,shapes.misc,shapes.arrows,intersections,decorations.pathmorphing,decorations.pathreplacing,backgrounds,shadows,babel}
\pgfplotsset{compat=1.18}
\usepgfplotslibrary{fillbetween}
\usepackage[np]{numprint}
\usepackage{listings,makecell,pas-tableur,twoopt}
\let\barre\undefined
\usepackage[pstricks]{tablvar}
\definecolor{darkgreen}{RGB}{0,100,0}
\providecommand{\red}{\color{red}}
\providecommand{\blue}{\color{blue}}
\providecommand{\green}{\color{green}}
\providecommand{\black}{\color{black}}
\providecommand{\white}{\color{white}}
\providecommand{\gray}{\color{gray}}
\providecommand{\cyan}{\color{cyan}}
\providecommand{\magenta}{\color{magenta}}
\providecommand{\yellow}{\color{yellow}}
\providecommand{\violet}{\color{violet}}
\providecommand{\orange}{\color{orange}}
\providecommand{\euro}{\texteuro}
\providecommand{\EUR}{\texteuro}
\providecommand{\ieme}{\textsuperscript{e}}
\providecommand{\degres}{\ensuremath{^\circ}}
\providecommand{\degree}{\ensuremath{^\circ}}
\providecommand{\vv}[1]{\overrightarrow{#1}}
\providecommand{\Coord}[1]{#1}
\providecommand{\decosix}{}
\providecommand{\starredbullet}{\ensuremath{\bullet}}
\pagestyle{empty}
\setlength{\parindent}{0pt}
\usepackage{adjustbox}\DeclareUnicodeCharacter{2236}{:}\DeclareUnicodeCharacter{2190}{\ensuremath{\leftarrow}}
\providecommand{\R}{\mathbb{R}}
\providecommand{\N}{\mathbb{N}}
\providecommand{\Z}{\mathbb{Z}}
\providecommand{\Q}{\mathbb{Q}}
\providecommand{\D}{\mathbb{D}}
\providecommand{\C}{\mathbb{C}}
\providecommand{\e}{{\mathrm{e}}}
\providecommand{\vect}[1]{\overrightarrow{#1}}
\providecommand{\vectt}[1]{\overrightarrow{\mathrm{#1}}}
\providecommand{\vvt}[1]{\overrightarrow{\mathrm{#1}}}
\providecommand{\np}[2][]{{#2~#1}}
\providecommand{\vv}[1]{{\overrightarrow{#1}}}
\renewcommand{\d}{{\mathrm{d}}}
\providecommand{\degres}{{^\circ}}
\providecommand{\degre}{{^\circ}}
\providecommand{\blue}{\color{blue}}
\providecommand{\red}{\color{red}}
\providecommand{\enskip}{\quad}
\providecommand{\llbracket}{\lbrack\!\lbrack}
\providecommand{\rrbracket}{\rbrack\!\rbrack}
\providecommand{\euro}{{\text{€}}}
\providecommand{\no}{{\text{n°}}}
\providecommand{\interval}[2]{{[#1;#2]}}
\providecommand{\intervFO}[2]{{[#1;#2[}}
\providecommand{\dotfill}{\ldots}

\providecommand{\ds}{\displaystyle}
\providecommand{\textding}[1]{#1}
\providecommand{\Oij}{\(\mathrm{O};\vec{\imath},\vec{\jmath}\)}
\providecommand{\Oijk}{\(\mathrm{O};\vec{\imath},\vec{\jmath},\vec{k}\)}
\providecommand{\Ouv}{\(\mathrm{O};\vec{u},\vec{v}\)}
\providecommand{\intervFF}[2]{\ensuremath{[#1;#2]}}
\renewcommand{\euro}{\ifmmode\text{\texteuro}\else\texteuro\fi}
\DeclareUnicodeCharacter{00B0}{\ifmmode^{\circ}\else\textdegree\fi}
\DeclareUnicodeCharacter{20AC}{\ifmmode\text{\texteuro}\else\texteuro\fi}
% Énoncés : Éducation nationale. Transcription et corrigés : APMEP, Denis Vergès et contributeurs. Découpage et classement : CoopMaths, Sébastien Lozano, Éric Elter et contributeurs. Adaptation de la présentation : Maths Explorer.
% CC BY-SA : https://coopmaths.fr/www/about/
% https://coopmaths.fr/alea/static/bac/2022/tex/bac_2022_05_sujet1_metropole_4_cor.tex
\begin{document}
\noindent\textit{APMEP / CoopMaths. CC BY-SA. Présentation adaptée par Maths Explorer.}\par\medskip
\medskip

\begin{enumerate}
	\item Soit $f$ définie sur $\R$ par $f(x)=\dfrac{-2x^2+3x-1}{x^2+1}$.\\
	Pour $x\neq 0$, $f(x)=\dfrac{x^2\left(-2+\dfrac{3}{x}-\dfrac{1}{x^2}\right)}{x^2\left(1+\dfrac{1}{x^2}\right)}=\dfrac{-2+\dfrac{3}{x}-\dfrac{1}{x^2}}{1+\dfrac{1}{x^2}}$ donc $\displaystyle\lim\limits_{x \rightarrow +\infty} f(x)=-2$.\\ Donc la droite horizontale d'équation $y=-2$ est une asymptote à la courbe représentative de la fonction $f$. \hfill \fbox{\textbf{Réponse c}}
	
	\item La fonction $f$ est continue sur $\R$ donc admet des primitives. Soit $F$ une primitive de $f$.
	
$\forall x \in \R$, $f(x)$ est de la forme $u'(x)\times \e^{u(x)}$, et admet pour primitives les fonctions $x\mapsto\e^{u(x)}+k$, $k\in\R$.

En posant $u(x) = x^2$ et $u'(x)=2x$, et en remarquant que $f(x)=\dfrac{1}{2}\times 2x~\e^{x^2}$, on peut donc écrire :

$\forall x \in \R$, $F(x)=\dfrac{1}{2}\times \e^{x^2}+k$.
	
	Sachant que $F(0)=1$ on en déduit que : $\dfrac{1}{2}+k=1$ donc $k=\dfrac{1}{2}$.

	Donc $\forall x \in \R$, $F(x)=\dfrac{1}{2}~\e^{x^2}+\dfrac{1}{2}$  \hfill \fbox{\textbf{Réponse d}}
	\item La courbe représentative est celle de la fonction $f'$. Avec la précision permise par le graphique, on peut affirmer que la fonction $f'$ est croissante sur $]-\infty~;~3]$ et décroissante sur $[3~;~+\infty[$. $[0~;~2] \subset ]-\infty~;~3]$, donc $f'$ est croissante $[0~;~2]$.\\
La fonction est donc convexe sur $[0~;~2]$. \hfill \fbox{\textbf{Réponse c}}
	\item La fonction $f$ définie sur $\R$ par $f(x)=3\e^{-x^2}+2$ est continue et dérivable. $\forall x \in \R$, $f(x)>0$.
	
Les primitives de la fonction $f$, ont pour dérivée $f$, qui est positive sur $\R$.

Donc les fonctions $F$ sont croissantes sur $\R$. \hfill \fbox{\textbf{Réponse a}}
	\item $\forall x \in ]0~;~+ \infty[$, $f(x)=\dfrac{2\ln(x)}{3x^2+1}= \dfrac23\dfrac{x^2\times \dfrac{\ln(x)}{x^2}}{x^2\left(1+\dfrac{1}{3x^2}\right)}= \dfrac23~\dfrac{\dfrac{\ln(x)}{x^2}}{1+\dfrac{1}{3x^2}}$.

Nous savons que : $\displaystyle\lim\limits_{x \rightarrow +\infty} \dfrac{\ln(x)}{x^2}= 0$ (croissances comparées) et
$\displaystyle\lim\limits_{x \rightarrow +\infty} 1+ \dfrac{1}{3x^2}= 1$

donc $\displaystyle\lim\limits_{x \rightarrow +\infty} f(x)= 0$ \hfill \fbox{\textbf{Réponse d}}
\item $\forall x \in \R$,  $(E)~~\e^{2x}+\e^x-12=0 \iff \left(\e^x\right)^2+\e^x-12=0$.

On pose comme changement de variable : $X=\e^x$ : $(E) \iff X^2+X-12=0$. Cette équation a pour solutions $X=-4$ et $X= 3$.

Or $X=\e^x>0$ donc l'unique solution sera celle de l'équation $X=\e^x=3\iff x=\ln(3)$. \hfill \fbox{\textbf{Réponse c}}
\end{enumerate}

\end{document}
